\chapter{Conformal Maps and the Open-String Scattering Prescription}

\lectureoverview{Conformal mapping first appears as a practical symmetry of
two-dimensional electrostatics.  Complex differentiability turns that
symmetry into a calculable one: the Cauchy--Riemann equations imply Laplace's
equation, and a direction-independent complex derivative preserves angles.
With this machinery in hand, a string worldsheet may be drawn as a strip,
half-plane, cylinder, or disc without changing the physics.  Open-string
scattering then becomes a Gaussian path integral on the disc with momentum
insertions on its boundary, followed by an integral over their inequivalent
positions.}

\section{Electrostatics in a Two-Dimensional World}

We begin with a literal two-dimensional electrostatic problem.  Charges and
their field lines are confined to a plane; they are not merely ordinary
three-dimensional charges whose motion happens to be planar.  Gauss's law
then gives
\begin{equation}
  F(r)\propto \frac{1}{r},
  \qquad
  V(r)\propto \log r,
  \label{eq:l08-two-dimensional-coulomb}
\end{equation}
because the flux crosses a circle whose circumference grows like \(r\).

Write the field in terms of a scalar potential \(\phi\).  In the sign
convention used on the board,
\begin{equation}
  \mathbf E=\nabla\phi,
  \qquad
  \nabla\cdot\mathbf E=\rho,
\end{equation}
and therefore
\begin{equation}
  \nabla^2\phi
  =\frac{\partial^2\phi}{\partial x^2}
  +\frac{\partial^2\phi}{\partial y^2}
  =\rho.
  \label{eq:l08-poisson}
\end{equation}
In a charge-free region this reduces to Laplace's equation,
\begin{equation}
  \nabla^2\phi=0.
  \label{eq:l08-laplace-phi}
\end{equation}
\editorialnote{Many electrodynamics texts define
\(\mathbf E=-\nabla\Phi\), in which case
\(\nabla^2\Phi=-\rho\) after constants are absorbed.  The lecture uses the
opposite potential sign.  The source-free equation and all conformal
arguments are unchanged.}

\begin{figure}[H]
  \centering
  \includegraphics[width=0.72\textwidth]{lecture_08_figure_01.png}
  \caption{The electrostatic equations unfolded into the
  two-dimensional Poisson operator.}
  \lectureframe{00:03:48}
  \label{fig:l08-electrostatics-board}
\end{figure}

The plus sign between the two spatial derivatives matters.  It makes this
an elliptic equation rather than a wave equation.  Conformal transformations
preserve its source-free form, so a solution on one domain may be carried to
a differently shaped domain together with its boundary conditions.  This is
useful in electrostatics, two-dimensional fluid flow, and cartography before
it is useful in string theory.

\begin{classroomqa}
\audiencequestion{If the Laplacian vanishes, is this a standing-wave
equation? \lecturetimestamp{00:05:35}}
\lecturerresponse{No.  Both spatial derivatives have the same sign and
there is no time derivative.  A wave equation requires time dependence with
the Lorentzian relative sign.  Electrostatics is static precisely because
nothing, including a wave, is evolving.}
\end{classroomqa}

\begin{classroomqa}
\audiencequestion{Is there a curl in two dimensions?
\lecturetimestamp{00:06:06}}
\lecturerresponse{Yes, but it has only one independent component.  For
\(\mathbf E=(E_x,E_y)\),
\[
  \operatorname{curl}_{2d}\mathbf E
  =\partial_yE_x-\partial_xE_y.
\]
It may be pictured as the component perpendicular to the board, although an
intrinsically two-dimensional world needs no external perpendicular
direction.  In higher dimensions the corresponding object is naturally an
antisymmetric two-form.}
\end{classroomqa}

\section{A Derivative with No Preferred Direction}

Introduce complex coordinates on the original and image planes:
\begin{equation}
  z=x+iy,
  \qquad
  w=u+iv,
  \qquad
  w=w(z).
  \label{eq:l08-complex-coordinates}
\end{equation}
Here \(x,y,u,v\) are real.  Locally, the derivative ought to be
\begin{equation}
  \frac{dw}{dz}
  =\lim_{\Delta z\to0}\frac{\Delta w}{\Delta z}.
  \label{eq:l08-complex-derivative}
\end{equation}
The complex plane, unlike the real line, permits infinitely many directions
of approach.  A continuous and ordinarily smooth map need not give the same
ratio along each direction.

Compare two simple approaches.  Along the \(x\)-direction, \(dy=0\), so
\begin{equation}
  \left.\frac{dw}{dz}\right|_x
  =\frac{\partial u}{\partial x}
  +i\frac{\partial v}{\partial x}.
\end{equation}
Along the \(y\)-direction, \(dx=0\), and
\begin{equation}
  \left.\frac{dw}{dz}\right|_y
  =\frac{1}{i}\frac{\partial u}{\partial y}
  +\frac{\partial v}{\partial y}
  =\frac{\partial v}{\partial y}
  -i\frac{\partial u}{\partial y}.
\end{equation}
Equality of real and imaginary parts gives
\begin{equation}
  \boxed{
  \frac{\partial u}{\partial x}
  =\frac{\partial v}{\partial y},
  \qquad
  \frac{\partial u}{\partial y}
  =-\frac{\partial v}{\partial x}}
  \label{eq:l08-cauchy-riemann}
\end{equation}
These are the Cauchy--Riemann equations.  Under the usual differentiability
assumptions they are also sufficient for the limit
\eqref{eq:l08-complex-derivative} to be direction independent.  A function
with this property is analytic, or holomorphic, at the point.

\begin{figure}[H]
  \centering
  \includegraphics[width=0.70\textwidth]{lecture_08_figure_02.png}
  \caption{The completed Cauchy--Riemann pair, including the relative
  minus sign.}
  \lectureframe{00:19:48}
  \label{fig:l08-cr-board}
\end{figure}

The lecture briefly distinguishes related terminology.  A function
holomorphic throughout the finite complex plane is entire.  A meromorphic
function may have isolated poles.  Thus \(1/z\) is analytic away from
\(z=0\), but it is not entire.

\section{Harmonic Conjugates and Conformal Geometry}

The Cauchy--Riemann equations are coupled first-order equations, but they
imply separate second-order equations.  Differentiate
\(\partial_xu=\partial_yv\) with respect to \(x\), differentiate
\(\partial_yu=-\partial_xv\) with respect to \(y\), and add:
\begin{align}
  \partial_x^2u&=\partial_x\partial_yv,\\
  \partial_y^2u&=-\partial_y\partial_xv,\\
  \therefore\qquad
  \partial_x^2u+\partial_y^2u&=0.
  \label{eq:l08-harmonic-u}
\end{align}
Repeating the argument with \(u\) and \(v\) interchanged gives
\begin{equation}
  \partial_x^2v+\partial_y^2v=0.
  \label{eq:l08-harmonic-v}
\end{equation}
The real and imaginary parts of an analytic function are therefore harmonic
conjugates.  Two arbitrary harmonic functions do not automatically make an
analytic function; they must still be linked by
\eqref{eq:l08-cauchy-riemann}.

\begin{figure}[H]
  \centering
  \includegraphics[width=0.68\textwidth]{lecture_08_figure_03.png}
  \caption{The mixed derivatives cancel, leaving Laplace's equation for
  one component of the analytic map.}
  \lectureframe{00:23:09}
  \label{fig:l08-harmonic-board}
\end{figure}

There is also a useful vector interpretation.  If
\(\mathbf F=(u,-v)\), then the two Cauchy--Riemann equations state, up to
the chosen orientation convention, that both the divergence and the
two-dimensional curl of \(\mathbf F\) vanish.  This is the promised return
to electrostatics.

\subsection{Why Angles Are Preserved}

Write two small complex displacements in polar form:
\begin{equation}
  \delta z=\rho e^{i\theta},
  \qquad
  \Delta z=\rho'e^{i\theta'}.
\end{equation}
Their ratio has phase \(\theta-\theta'\), exactly the angle between them.
Under an analytic map,
\begin{equation}
  \frac{\delta w}{\delta z}
  =\frac{\Delta w}{\Delta z}
  =w'(z),
\end{equation}
because the derivative does not depend on direction.  Rearranging gives
\begin{equation}
  \frac{\Delta z}{\delta z}
  =\frac{\Delta w}{\delta w}.
  \label{eq:l08-angle-ratio}
\end{equation}
The two sides have the same phase, so the angle between the original
intervals equals the angle between their images.

\begin{figure}[H]
  \centering
  \includegraphics[width=0.70\textwidth]{lecture_08_figure_04.png}
  \medskip

  \begin{tikzpicture}[x=0.9cm,y=0.9cm,>=Latex]
    \begin{scope}
      \draw[->] (-0.2,0) -- (2.3,0);
      \draw[->] (0,-0.2) -- (0,2.1);
      \draw[very thick,blue!65!black,->] (0,0) -- (1.85,0.45)
        node[right] {\(\delta z\)};
      \draw[very thick,orange!75!black,->] (0,0) -- (0.72,1.65)
        node[above] {\(\Delta z\)};
      \draw (0.58,0.14) arc[start angle=14,end angle=66,radius=0.60];
      \node at (1.0,-0.55) {\(z\)-plane};
    \end{scope}
    \draw[->,thick] (3.0,0.9) -- (4.2,0.9)
      node[midway,above] {\(w(z)\)};
    \begin{scope}[xshift=5.0cm]
      \draw[->] (-0.2,0) -- (2.3,0);
      \draw[->] (0,-0.2) -- (0,2.1);
      \draw[very thick,blue!65!black,->] (0,0) -- (1.65,0.70)
        node[right] {\(\delta w\)};
      \draw[very thick,orange!75!black,->] (0,0) -- (0.48,1.62)
        node[above] {\(\Delta w\)};
      \draw (0.54,0.23) arc[start angle=23,end angle=74,radius=0.59];
      \node at (1.0,-0.55) {\(w\)-plane};
    \end{scope}
  \end{tikzpicture}
  \caption{The board proof and a clean local reconstruction.  Multiplication
  by \(w'(z)\) rotates and rescales every infinitesimal direction by the
  same amount, preserving their relative angle.}
  \lectureframe{00:34:11}
  \label{fig:l08-angle-proof}
\end{figure}

This local argument assumes \(w'(z)\neq0\).  At a critical point the map may
fail to be locally one-to-one, and one must treat the geometry separately.
Away from such points, analytic maps and orientation-preserving conformal
maps are locally the same objects.

\section{Examples that Train the Eye}

\subsection{Squaring and Complex Conjugation}

For \(w=z^2\),
\begin{equation}
  (x+iy)^2=x^2-y^2+2ixy=u+iv,
\end{equation}
so
\begin{equation}
  u=x^2-y^2,
  \qquad
  v=2xy.
  \label{eq:l08-z-squared-components}
\end{equation}
The second derivatives of \(u\) are \(2\) and \(-2\); those of \(v\)
vanish separately.  Both are harmonic, and direct differentiation also
checks the Cauchy--Riemann equations.

\begin{figure}[H]
  \centering
  \includegraphics[width=0.72\textwidth]{lecture_08_figure_05.png}
  \caption{The squaring map separated into its real and imaginary
  harmonic components.}
  \lectureframe{00:36:45}
  \label{fig:l08-z-square-board}
\end{figure}

\begin{classroomqa}
\audiencequestion{Does \(1/z\) give a nonanalytic example?
\lecturetimestamp{00:37:49}}
\lecturerresponse{Only at one point.  It is analytic wherever \(z\neq0\)
and has a pole at the origin.  To see a map that fails the
Cauchy--Riemann test everywhere, complex conjugation is a better example.}
\end{classroomqa}

By contrast,
\begin{equation}
  w=\bar z=x-iy
  \qquad\Longrightarrow\qquad
  u=x,\quad v=-y
\end{equation}
reverses orientation.  Its derivatives fail the first
Cauchy--Riemann equation by a sign, so it is not analytic as a function of
\(z\).

\begin{figure}[H]
  \centering
  \includegraphics[width=0.62\textwidth]{lecture_08_figure_06.png}
  \caption{Complex conjugation provides the deliberately nonanalytic
  counterexample.}
  \lectureframe{00:38:38}
  \label{fig:l08-conjugate-board}
\end{figure}

\subsection{The Exponential and Its Inverse}

\begin{classroomqa}
\audiencequestion{If \(e^z\) is analytic, why should its inverse
\(\log z\) also be analytic? \lecturetimestamp{00:40:38}}
\lecturerresponse{Where the inverse exists and \(w'(z)\neq0\), its
derivative is \(1/w'(z)\).  Direction independence is therefore preserved.
Globally the logarithm is multivalued, so one chooses a branch and excludes
its branch point and cut.  The local conformal statement remains valid on
that chosen domain.}
\end{classroomqa}

Euler's formula gives
\begin{equation}
  w=e^z=e^{x+iy}
  =e^x(\cos y+i\sin y),
\end{equation}
hence
\begin{equation}
  u=e^x\cos y,
  \qquad
  v=e^x\sin y.
  \label{eq:l08-exponential-components}
\end{equation}
For example,
\(\partial_xu=e^x\cos y=\partial_yv\), and the second
Cauchy--Riemann equation follows just as directly.

\begin{figure}[H]
  \centering
  \includegraphics[width=0.68\textwidth]{lecture_08_figure_07.png}
  \caption{The real and imaginary components of the exponential map.}
  \lectureframe{00:42:23}
  \label{fig:l08-exponential-board}
\end{figure}

\section{Half-Plane, Strip, Cylinder, and Disc}

\subsection{The Logarithm Turns Radius into Time}

Take the right half-plane, where
\(-\pi/2<\theta<\pi/2\), and write
\begin{equation}
  z=re^{i\theta},
  \qquad
  w=\log z=\log r+i\theta.
  \label{eq:l08-log-map}
\end{equation}
The half-plane maps to the infinite strip
\begin{equation}
  -\frac{\pi}{2}<\operatorname{Im}w<\frac{\pi}{2}.
\end{equation}
Rays of fixed \(\theta\) become horizontal lines, while semicircles of fixed
\(r\) become vertical lines.  The origin is sent to
\(\operatorname{Re}w=-\infty\), and \(r=\infty\) to
\(\operatorname{Re}w=+\infty\).

\begin{figure}[H]
  \centering
  \includegraphics[width=0.75\textwidth]{lecture_08_figure_08.png}
  \caption{The logarithm carries the right half-plane into an infinite
  strip: polar rays and semicircles become an orthogonal rectangular grid.}
  \lectureframe{00:47:48}
  \label{fig:l08-log-strip-board}
\end{figure}

The strip is the standard worldsheet of an open string.  Its bounded
coordinate is the parameter along the string, while translation along the
unbounded direction is worldsheet time.  Under the inverse exponential map,
a time translation becomes a dilation \(z\mapsto e^a z\).  Thus the same
symmetry looks like translation on the strip and uniform expansion on the
half-plane.

\begin{classroomqa}
\audiencequestion{Where does a particular value such as \(r=1\) go, and
are the coordinates on the strip still polar? \lecturetimestamp{00:50:58}}
\lecturerresponse{A circle of fixed radius becomes a vertical line at
\(\operatorname{Re}w=\log r\); in particular \(r=1\) maps to the central
line \(\operatorname{Re}w=0\).  Polar coordinates describe the
\(z\)-plane.  The \(w\)-plane uses ordinary rectangular coordinates
\((\log r,\theta)\).}
\end{classroomqa}

If the full punctured plane is mapped by \(\log z\), the angular coordinate
has period \(2\pi\).  The resulting strip therefore has its upper and lower
edges identified: it is an infinite cylinder.  The half-plane with a
boundary describes the open-string worldsheet; the full plane, with no
boundary and periodic angle, describes a closed-string cylinder.

\begin{figure}[H]
  \centering
  \includegraphics[width=0.75\textwidth]{lecture_08_figure_09.png}
  \caption{The half-plane/strip picture and the periodic identification
  that turns the full punctured plane into a closed-string cylinder.}
  \lectureframe{00:54:23}
  \label{fig:l08-open-closed-board}
\end{figure}

\subsection{A Cayley Map to the Unit Disc}

A linear fractional, or Möbius, map carries the right half-plane to the unit
disc:
\begin{equation}
  w=\frac{z-1}{z+1}.
  \label{eq:l08-cayley}
\end{equation}
It sends \(z=0\) to \(-1\), \(z=\infty\) to \(+1\), and the imaginary axis
to \(|w|=1\).  If \(z=ir\),
\begin{equation}
  \left|\frac{ir-1}{ir+1}\right|=1.
\end{equation}
Points with \(\operatorname{Re}z>0\) lie inside the circle; the left
half-plane maps outside it.
\editorialnote{During the live derivation the reciprocal
\((z+1)/(z-1)\) is written while the interior-disc image is described.
That reciprocal has the same unit-circle boundary but sends the right
half-plane to the exterior.  Equation \eqref{eq:l08-cayley} is the
interior-disc convention consistent with the stated geometry.}

\begin{figure}[H]
  \centering
  \includegraphics[width=0.72\textwidth]{lecture_08_figure_10.png}
  \caption{The linear-fractional construction of the unit-circle boundary.
  The reciprocal convention visible on the board exchanges interior and
  exterior.}
  \lectureframe{00:58:38}
  \label{fig:l08-cayley-board}
\end{figure}

Möbius maps send the generalized collection of circles and straight lines
to circles and straight lines.  The radial and circular grid on the
half-plane consequently becomes two orthogonal families of arcs on the
disc.  A straight diameter is simply the limiting case of a circle whose
center has gone to infinity.

\begin{classroomqa}
\audiencequestion{Are the arcs drawn inside the disc really circles?
\lecturetimestamp{01:00:55}}
\lecturerresponse{Yes.  Their centers occur at different locations, often
outside the visible disc.  The diameter is the infinite-radius limit of the
same family.}
\end{classroomqa}

\begin{classroomqa}
\audiencequestion{Does the disc drawing mean that it is a sphere, and where
does the other half-plane go? \lecturetimestamp{01:01:21}}
\lecturerresponse{The drawing here is a flat disc, not a sphere.  The right
half-plane maps inside its boundary and the left half-plane maps outside.
If one adds the point at infinity to the entire complex plane, that extended
plane can be represented as a sphere by stereographic projection; that is a
different global statement.}
\end{classroomqa}

\begin{classroomqa}
\audiencequestion{Which disc curves correspond to the original radial
spokes and constant-radius semicircles? \lecturetimestamp{01:02:59}}
\lecturerresponse{The positive real axis becomes the diameter joining
\(-1\) and \(+1\).  Other rays become circular arcs joining those same
boundary points.  Semicircles of fixed radius become the orthogonal family
that begins and ends on the unit circle.}
\end{classroomqa}

\begin{figure}[H]
  \centering
  \includegraphics[width=0.72\textwidth]{lecture_08_figure_11.png}
  \caption{Lines and semicircles in the half-plane become orthogonal
  circle-and-line families in the disc.}
  \lectureframe{01:03:48}
  \label{fig:l08-circle-grid-board}
\end{figure}

\section{The Four-String Worldsheet}

We can now return to scattering.  Two open strings enter, join into one
string, propagate for an arbitrary interval, and split into two outgoing
strings.  In strip coordinates the worldsheet looks like a long ribbon slit
from each end---the "sports bandage" picture.  The intermediate lifetime
is not chosen in advance.  It is one of the histories over which the
amplitude is integrated.

Let \((\sigma,\tau)\) be coordinates on the worldsheet and let
\begin{equation}
  X^\mu=X^\mu(\sigma,\tau)
  \label{eq:l08-embedding}
\end{equation}
give its position in spacetime.  The labels \(\sigma,\tau\) live on the
two-dimensional parameter surface; the index \(\mu\) labels spacetime
coordinates.  In the bosonic theory \(\mu=0,\ldots,25\); in the superstring
case there are ten spacetime coordinates.

After Wick rotation, the free embedding fields have Euclidean action
\begin{equation}
  S_E[X]
  =\frac{1}{4\pi\alpha'}\int_\Sigma d\sigma\,d\tau\,
  \left[
    \partial_\tau X^\mu\partial_\tau X_\mu
    +\partial_\sigma X^\mu\partial_\sigma X_\mu
  \right].
  \label{eq:l08-euclidean-action}
\end{equation}
The corresponding path integral sums over every embedding of the chosen
worldsheet into spacetime:
\begin{equation}
  \int\mathcal D X\,e^{-S_E[X]}.
  \label{eq:l08-bare-worldsheet-integral}
\end{equation}

\begin{figure}[H]
  \centering
  \includegraphics[width=0.74\textwidth]{lecture_08_figure_12.png}
  \caption{The completed positive Euclidean worldsheet weight.  The two
  derivative terms are summed over all spacetime components.}
  \lectureframe{01:11:34}
  \label{fig:l08-action-board}
\end{figure}

The slit strip has one connected boundary.  Conformal freedom therefore
allows it to be represented as a disc.  Four remote ends of the strips
become four marked points on the boundary.  Deforming the marked points
changes the apparent separation between joining and splitting in strip
coordinates, but it does not change the underlying surface or physics.

\begin{figure}[H]
  \centering
  \includegraphics[width=0.76\textwidth]{lecture_08_figure_13.png}
  \caption{The slit-strip worldsheet and its disc representation with
  boundary points for the external strings.}
  \lectureframe{01:15:18}
  \label{fig:l08-bandage-disc-board}
\end{figure}

The detailed grid lines used to picture this conformal map carry no physical
content.  What matters is the topology, the boundary, and the marked
external points.  The same worldsheet can be stretched into different
channel pictures without introducing a new interaction.

\section{External States as Boundary Insertions}

The bare integral \eqref{eq:l08-bare-worldsheet-integral} does not yet say
which momenta enter and leave.  For a scalar open-string state of momentum
\(k_i^\mu\), insert on the boundary the vertex operator
\begin{equation}
  V_i(z_i)=e^{\,i k_i\cdot X(z_i)},
  \qquad
  k_i\cdot X=k_{i\mu}X^\mu.
  \label{eq:l08-vertex}
\end{equation}
The point \(z_i\) is a position on the worldsheet boundary.  It is not a
spacetime coordinate and not a direction of momentum.

\begin{figure}[H]
  \centering
  \includegraphics[width=0.72\textwidth]{lecture_08_figure_14.png}
  \caption{The bulk Euclidean weight and the product of boundary momentum
  insertions are displayed as distinct factors.}
  \lectureframe{01:18:58}
  \label{fig:l08-vertex-board}
\end{figure}

For \(n\) external open strings, the schematic prescription is
\begin{equation}
  \begin{aligned}
    \mathcal A_n(k_1,\ldots,k_n)
    &=\int_{\mathcal M_{0,n}^{\partial}}\!d\mu(z_i)
      \int\mathcal D X\\
    &\quad\times e^{-S_E[X]}
      \prod_{i=1}^{n}e^{\,i k_i\cdot X(z_i)}.
  \end{aligned}
  \label{eq:l08-open-string-prescription}
\end{equation}
The functional integral is over embeddings \(X^\mu\).  The finite-dimensional
integral is over inequivalent boundary insertion positions.  Polarization
factors and the gauge-fixing measure are suppressed at this stage, as they
are in the lecture.
\editorialnote{A complete gauge-fixed string amplitude also contains the
appropriate ghost contribution, state-dependent vertex factors, coupling
constants, and normalization.  Equation
\eqref{eq:l08-open-string-prescription} records the embedding-field and
position dependence developed in this lecture.}

\begin{classroomqa}
\audiencequestion{If the external momenta are fixed, why are they integrated
over? \lecturetimestamp{01:20:42}}
\lecturerresponse{They are not.  The \(k_i\) are fixed labels of the
scattering amplitude.  What is integrated are the embedding fields and the
worldsheet locations \(z_i\) at which those fixed momenta are inserted.
The latter integral replaces the integral over the intermediate lifetime in
the strip picture.}
\end{classroomqa}

\begin{classroomqa}
\audiencequestion{Is the product of vertex factors inside the path integral,
and is it part of the action? \lecturetimestamp{01:21:39}}
\lecturerresponse{It multiplies \(e^{-S_E}\) inside the functional
integral.  One may equivalently exponentiate the factors and regard them as
localized source terms added to the action, but they are conceptually vertex
operators rather than part of the free bulk action.}
\end{classroomqa}

\begin{figure}[H]
  \centering
  \includegraphics[width=0.76\textwidth]{lecture_08_figure_15.png}
  \caption{The path integral written with separate factors
  \(e^{ik_i\cdot X(z_i)}\) for successive external particles.}
  \lectureframe{01:23:23}
  \label{fig:l08-full-prescription-board}
\end{figure}

\begin{classroomqa}
\audiencequestion{Why does \(e^{ik\cdot X}\) still contain \(i\) after the
worldsheet Wick rotation? \lecturetimestamp{01:23:42}}
\lecturerresponse{The rotation changes the worldsheet time and converts the
bulk weight from \(e^{iS}\) to \(e^{-S_E}\).  The \(i\) in a momentum
eigenfunction is the Fourier phase that labels spacetime momentum.  It has a
different role and remains.}
\end{classroomqa}

\begin{classroomqa}
\audiencequestion{Does the position \(z_i\) determine the direction of the
momentum vector? \lecturetimestamp{01:24:19}}
\lecturerresponse{No.  Direction is encoded in the spacetime components
\(k_i^\mu\).  The insertion point \(z_i\) says where that momentum enters
the worldsheet description; \(X^\mu(z_i)\) is among the field values summed
over by \(\mathcal D X\).}
\end{classroomqa}

\begin{classroomqa}
\audiencequestion{Do the drawn angles of the bandage legs encode the
directions of the physical momenta? \lecturetimestamp{01:25:07}}
\lecturerresponse{No.  The bandage is parameter space, not an image of the
string in physical spacetime.  Its legs may be drawn horizontally, tilted,
or conformally bent.  Physical directions are carried by the components
\(k^\mu\), while \(X^\mu(\sigma,\tau)\) is integrated over.}
\end{classroomqa}

\section{The Gaussian Integral as Electrostatics}

The action \eqref{eq:l08-euclidean-action} is quadratic in every component
\(X^\mu\), and the insertions are linear sources once their exponentials are
combined with the action.  The functional integral is therefore Gaussian.
At its saddle point, each spacetime component obeys a Poisson equation of
the schematic form
\begin{equation}
  \nabla^2 X^\mu(z)
  \propto
  i\alpha'\sum_i k_i^\mu\,\delta(z-z_i),
  \label{eq:l08-electrostatic-source}
\end{equation}
with the open-string boundary condition on the rim of the disc.  Each
\(X^\mu\) behaves like an independent electrostatic potential, and each
momentum component \(k_i^\mu\) behaves like the corresponding charge placed
at \(z_i\).

\begin{figure}[H]
  \centering
  \includegraphics[width=0.76\textwidth]{lecture_08_figure_16.png}
  \caption{The disc with boundary sources is interpreted as one independent
  two-dimensional electrostatic problem for each spacetime component.}
  \lectureframe{01:28:58}
  \label{fig:l08-electrostatic-disc-board}
\end{figure}

Solve the electrostatic problem, evaluate its energy, exponentiate the
negative energy, and finally integrate over the inequivalent charge
positions.  In standard open-string normalization, the embedding-field
correlator supplies the Koba--Nielsen factor
\begin{equation}
  \left\langle\prod_i e^{\,ik_i\cdot X(z_i)}\right\rangle
  \propto
  (2\pi)^D\delta^{(D)}\!\left(\sum_i k_i\right)
  \prod_{i<j}|z_i-z_j|^{2\alpha' k_i\cdot k_j}.
  \label{eq:l08-koba-nielsen}
\end{equation}
The delta function expresses momentum conservation.  The product contains
the electrostatic pair energies.  For the four scalar boundary insertions
used here, fixing three positions turns the remaining integral into the
Euler beta integral encountered in the preceding lectures.

\begin{classroomqa}
\audiencequestion{Are there twenty-six electrostatic energies or only one?
\lecturetimestamp{01:30:53}}
\lecturerresponse{There is one independent potential and field for each
spacetime component.  Their quadratic energies have no cross terms in an
orthonormal target-space basis, and the contributions are added to one total
action before exponentiation.}
\end{classroomqa}

\begin{classroomqa}
\audiencequestion{Is the angular coordinate of the disc the physical
scattering angle? \lecturetimestamp{01:32:10}}
\lecturerresponse{No.  It is a worldsheet position parameter.  Its
inequivalent combination is related to the intermediate propagation
interval and is integrated.  The physical scattering angle is encoded
through Lorentz invariants built from the fixed momenta.}
\end{classroomqa}

\begin{classroomqa}
\audiencequestion{Why does this construction fail away from the critical
dimension? \lecturetimestamp{01:32:35}}
\lecturerresponse{The formal Gaussian integrals alone are not the whole
consistency condition.  Gauge fixing introduces a conformal anomaly whose
matter and ghost contributions cancel only at the critical dimension:
\(D=26\) for the bosonic string and \(D=10\) for the superstring.  Without
that cancellation the supposed conformal symmetry of the quantum
worldsheet is lost.}
\end{classroomqa}

\begin{classroomqa}
\audiencequestion{Have we already assumed three large spatial dimensions
and compact extra dimensions? \lecturetimestamp{01:33:38}}
\lecturerresponse{No.  Compactification has not yet entered the
construction.  At this stage the bosonic coordinates still range over all
twenty-six target-space directions, or ten in the superstring case.  How a
subset can remain macroscopically large is the subject of the next
discussion.}
\end{classroomqa}

\begin{classroomqa}
\audiencequestion{Is the connection with electrostatics simply that the
coordinates satisfy Laplace's equation? \lecturetimestamp{01:33:57}}
\lecturerresponse{Exactly.  Away from the insertion points each
\(X^\mu\) is harmonic; the vertex operators supply delta-function sources.
That is mathematically the same boundary-value problem as a static
two-dimensional potential with point charges, whether or not we use the
electrostatic vocabulary.}
\end{classroomqa}

\subsection{Conformal Redundancy and the Number of Moduli}

The disc has a three-real-parameter conformal automorphism group,
\(\mathrm{PSL}(2,\mathbb R)\).  Three ordered boundary points may therefore
be fixed at convenient locations.  For four external strings only the
fourth position remains to be integrated; for \(n\) ordered boundary
insertions there are
\begin{equation}
  n-3
  \label{eq:l08-n-minus-three}
\end{equation}
real position integrals at tree level.

\begin{classroomqa}
\audiencequestion{Are there four independent angular positions around the
disc? \lecturetimestamp{01:34:37}}
\lecturerresponse{Not physically.  One initially writes four positions, but
disc conformal symmetry can place any three wherever convenient.  Only the
relative location of the fourth remains for the two-in/two-out amplitude.}
\end{classroomqa}

\begin{figure}[H]
  \centering
  \includegraphics[width=0.72\textwidth]{lecture_08_figure_17.png}
  \caption{Four boundary insertions on the disc.  Conformal freedom fixes
  any three, leaving one four-point modulus.}
  \lectureframe{01:35:18}
  \label{fig:l08-fix-three-board}
\end{figure}

\begin{classroomqa}
\audiencequestion{What if the two split points do not occur at the same
height in the strip drawing? \lecturetimestamp{01:35:24}}
\lecturerresponse{A conformal map can move that apparent mismatch away.
Different drawings are coordinate representatives of the same modulus.
The invariant integration is the one over the remaining boundary position,
equivalently over the propagation interval.}
\end{classroomqa}

\begin{classroomqa}
\audiencequestion{For more external particles, how many position parameters
remain? \lecturetimestamp{01:36:00}}
\lecturerresponse{Fix three insertion points with conformal symmetry and
integrate the other \(n-3\).  This is the boundary-moduli count for an
ordered tree-level open-string amplitude.}
\end{classroomqa}

\begin{classroomqa}
\audiencequestion{How is this related to M-theory?
\lecturetimestamp{01:36:37}}
\lecturerresponse{Not yet.  The construction here is perturbative
open-string scattering on a worldsheet.  Its relation to M-theory requires
later duality and nonperturbative ideas, so the question is explicitly
deferred.}
\end{classroomqa}

\begin{classroomqa}
\audiencequestion{Is a different vertical or crossed strip picture a
different calculation? \lecturetimestamp{01:36:49}}
\lecturerresponse{No.  One may hold a convenient point fixed and let the
remaining modulus range through its full interval.  The strip pictures and
the boundary-point integral on the disc are different slicings of the same
integral, which is why all external strings can be treated symmetrically.}
\end{classroomqa}

\section{What the Prescription Does Not Yet Solve}

The disc formula extends immediately to more external strings: add more
vertex operators and integrate their \(n-3\) inequivalent boundary
positions.  It is a precise tree-level prescription for a consistent string
model, but it is not yet a realistic model of observed particle physics.
The simple bosonic theory has the wrong spacetime dimension and additional
physical defects.  Superstrings improve the spectrum and lower the critical
dimension to ten, but six spatial dimensions still have to be compactified.

\begin{classroomqa}
\audiencequestion{Does this simple model already agree with particle
collision experiments? \lecturetimestamp{01:39:06}}
\lecturerresponse{No.  The resonance structure historically resembles
hadronic scattering, but the twenty-six-dimensional bosonic model cannot be
a direct description of observed particles.  The formalism defines
scattering amplitudes for the model; matching nature requires a realistic
spectrum, compactification, interactions, and vacuum choice.}
\end{classroomqa}

\begin{classroomqa}
\audiencequestion{Will compactification reduce the extra dimensions and
produce realistic particle data? \lecturetimestamp{01:41:10}}
\lecturerresponse{It can reduce the effective macroscopic dimension and
generate particle spectra and gauge groups.  It is also where much of the
real mathematical difficulty enters.  Many consistent compactifications
are known, and selecting a unique model of observed elementary particles
remains unresolved.}
\end{classroomqa}

The lecture emphasizes this distinction rather than advertising a finished
theory.  String constructions are mathematically definite and have produced
deep information about gravity and geometry, yet their multiplicity is
itself part of the physical problem.

\begin{classroomqa}
\audiencequestion{Is there a short way to remember why the bosonic critical
dimension is twenty-six? \lecturetimestamp{01:43:19}}
\lecturerresponse{The earlier oscillator calculation used the regularized
zero-point sum \(1+2+3+\cdots=-1/12\).  The transverse oscillator
contribution involves twenty-four directions, and restoring the two
longitudinal spacetime coordinates gives \(24+2=26\).  This is a mnemonic
for the earlier consistency argument, not a new derivation here.}
\end{classroomqa}

\begin{classroomqa}
\audiencequestion{Are the inserted momenta \(k_i\) themselves quantized?
\lecturetimestamp{01:43:46}}
\lecturerresponse{In noncompact spacetime their components are continuous.
What is quantized is the oscillator level and hence the allowed mass-squared
spectrum.  Compact dimensions may quantize the corresponding momentum and
winding components, but that belongs to the next discussion.}
\end{classroomqa}
